\section{MFA aktyvavimas}
\label{sec:aktyvavimas}

MFA gali aktyvuoti tik prisijungęs vartotojas. Aktyvavimo metu jo paskyra susiejama su \textit{Google} paskyra, t.~y. išsaugomas \textit{Google} ID žetono \texttt{sub}. Aktyvavimo seka pavaizduota \ref{fig:aktyvavimas} pav.:
\begin{enumerate}
    \item Vartotojas nustatymų skiltyje „Saugumas“ paspaudžia „Įjungti MFA per Google“.
    \item Aplikacija paprašo pakartotinai įvesti slaptažodį ir jį patikrina. Taip kitas asmuo negalės įjungti MFA pasinaudojęs palikta atvira sesija.
    \item Aplikacija sugeneruoja \texttt{state}, \texttt{nonce} ir \texttt{code\_verifier}. Jie išsaugomi serverio sesijoje kartu su tikslu \texttt{purpose = enroll} ir galiojimo laiku (10 min.).
    \item Naršyklė peradresuojama į \textit{Google} su \texttt{scope = openid email} ir \texttt{prompt = select\_account}. Taip vartotojas sąmoningai pasirenka, kurią \textit{Google} paskyrą susieti.
    \item Vartotojas prisijungia prie \textit{Google} ir sutikimo lange (\textit{angl. consent screen}) leidžia aplikacijai gauti jo identifikatorių ir el. pašto adresą.
    \item \textit{Google} grąžina vartotoją į \texttt{redirect\_uri}. Aplikacija patikrina \texttt{state}, iškeičia kodą į žetonus (su \texttt{code\_verifier}), patikrina ID žetoną ir \texttt{nonce} (žr.~\ref{subsec:srautas} poskyrį).
    \item Aplikacija patikrina, ar pora (\texttt{google}, \texttt{sub}) dar nėra susieta su kita paskyra. Jei susieta, aktyvavimas nutraukiamas.
    \item Duomenų bazėje išsaugomas susiejimas (\texttt{oauth\_identities} lentelė) ir nustatoma \texttt{users.mfa\_enabled = true}.
    \item Sugeneruojama 10 vienkartinių atkūrimo kodų. Jie parodomi vartotojui tik vieną kartą, o duomenų bazėje saugomos tik jų maišos (argon2id). Tai vienas iš skaidrėse siūlomų MFA atkūrimo būdų~\cite{VMA}.
    \item Vartotojui išsiunčiamas el. laiškas „MFA įjungtas“. Jei MFA įjungė ne jis, apie tai sužinos iš karto.
\end{enumerate}

\begin{figure}[H]
    \centering
    \begin{adjustbox}{max width=\textwidth}
    \begin{tikzpicture}
        % Dalyviai
        \node[actor] (U) at (0,0)  {Vartotojas\\(naršyklė)};
        \node[actor] (A) at (5,0)  {Aplikacija\\(serveris)};
        \node[actor] (G) at (10,0) {\textit{Google}\\(OpenID teikėjas)};
        \node[actor] (D) at (15,0) {Duomenų\\bazė};

        % Gyvavimo linijos
        \foreach \x in {0,5,10,15} {
            \draw[lifeline] (\x,-0.5) -- (\x,-21.8);
        }

        \draw[msg] (0,-1.4) -- node[lbl, above] {1. „Įjungti MFA per Google“} (5,-1.4);

        \draw[msg] (0,-2.4) -- node[lbl, above] {2. Slaptažodis (pakartotinai)} (5,-2.4);

        \draw[msg] (5,-3.4) -- node[lbl, above, pos=0.75] {2. Slaptažodžio patikra (argon2id)} (15,-3.4);

        \node[note] at (5,-4.6) {3. \texttt{state}, \texttt{nonce}, \texttt{code\_verifier} $\to$ sesija\\
            (\texttt{purpose = enroll}, galioja 10 min.)};

        \draw[ret] (5,-6.2) -- node[lbl, above] {4. HTTP 302 $\to$ \textit{Google}} (0,-6.2);

        \draw[msg] (0,-7.4) -- node[lbl, above, pos=0.25] {4. \texttt{scope=openid email},\\
            \texttt{prompt=select\_account}} (10,-7.4);

        \node[note] at (10,-8.5) {5. Prisijungimas prie \textit{Google}\\
            ir sutikimas (\textit{consent})};

        \draw[ret] (10,-10.2) -- node[lbl, above, pos=0.75] {6. HTTP 302 $\to$ \texttt{redirect\_uri}\\
            (\texttt{code}, \texttt{state})} (0,-10.2);

        \draw[msg] (0,-11.4) -- node[lbl, above] {6. GET \texttt{/auth/google/callback}} (5,-11.4);

        \draw[msg] (5,-12.6) -- node[lbl, above] {6. POST \texttt{/token} (+\texttt{code\_verifier})} (10,-12.6);

        \draw[ret] (10,-13.6) -- node[lbl, above] {6. \texttt{id\_token}} (5,-13.6);

        \node[note] at (5,-14.6) {6. Tikrina \texttt{state}, ID žetoną, \texttt{nonce}};

        \draw[msg] (5,-15.8) -- node[lbl, above, pos=0.75] {7. Ar (\texttt{google}, \texttt{sub}) jau susietas?} (15,-15.8);

        \draw[ret] (15,-16.8) -- node[lbl, above, pos=0.25] {7. Ne} (5,-16.8);

        \draw[msg] (5,-18.0) -- node[lbl, above, pos=0.75] {8. INSERT \texttt{oauth\_identities};\\
            \texttt{mfa\_enabled = true}} (15,-18.0);

        \draw[msg] (5,-19.2) -- node[lbl, above, pos=0.75] {9. INSERT 10 \texttt{code\_hash} (argon2id)} (15,-19.2);

        \draw[ret] (5,-20.2) -- node[lbl, above] {9. Atkūrimo kodai (rodomi 1 kartą)} (0,-20.2);

        \draw[msg] (5,-21.2) -- node[lbl, above] {10. El. laiškas „MFA įjungtas“} (0,-21.2);
    \end{tikzpicture}
    \end{adjustbox}
    \caption{MFA aktyvavimo seka.}
    \label{fig:aktyvavimas}
\end{figure}


MFA išjungimui ir susietos \textit{Google} paskyros pakeitimui reikalingas sėkmingas OAuth patvirtinimas (žr.~\ref{subsec:step_up} poskyrį). Kitaip užpuolikas, sužinojęs tik slaptažodį, galėtų MFA tiesiog išjungti.
